\pdfminorversion=7%
\documentclass[aspectratio=169,mathserif,notheorems]{beamer}%
%
\xdef\bookbaseDir{../../bookbase}%
\xdef\sharedDir{../../shared}%
\RequirePackage{\bookbaseDir/styles/slides}%
\RequirePackage{\sharedDir/styles/styles}%
\toggleToGerman%
%
\subtitle{27.~Konzeptuelles Schema:~Entitäten und Attribute}%
%
\begin{document}%
%
\startPresentation%
%
\section{Einleitung}%
%
\begin{frame}%
\frametitle{Einleitung}%
\begin{itemize}%
\item Beginnen wir nun mit dem konzeptuellen Modellieren von Applikationen.%
%
\item<2-> Durch die Anforderungsanalyse haben wir ein klares Verständnis über die Entitäten in unserem Szenario und die Beziehungen zwischen ihnen.%
%
\item<3-> In unserem Beispiel haben wir die Entitäten der Plattform für der Management der Lehre einer Universität kennengelernt.%
%
\item<4-> Wir haben sie aber nur sehr informell diskutiert.%
%
\item<5-> Jetzt wollen wir die Informationen über die Entitäten in ein konsistentes Modell überführen.%
%
\item<6-> Im konzeptuellen Designschritt sind diese Modelle unabhängig von konkreten Technologien.%
%
\item<7-> Dieser Schritt wird \emph{entity relationship modeling} genannt\cite{G2011EW2ITDS:CMUTERM,SS2005EIDDDFDB:CDDRAAML,SS2005EIDDDFDB:CDDICAMP,V1999C5DMS:CDUTERM,B1990CMERMO}.%
%
\item<8-> Er ist besonders bei großen Projekten sinnvoll.%
%
\item<9-> Bei kleinen Projekten, wie \DEzB\ einer \glslink{db}{Datenbank} für unsere eigenen Literaturreferenzen oder für unsere Musiksammlung kann man ihn überspringen und direkt zum logischen Design übergehen~\cite{S2024D:CDMERDE}.%
%
\item<10-> Unsere Lehrmanagementplattform ist aber ein großes Projekt und man muss daher mit dem konzeptuellen Design beginnen.%
\end{itemize}%
\end{frame}%
%
\section{Entitäten und Attribute}%
%
\begin{frame}%
\frametitle{Entitäten}%
%
\begin{itemize}%
\item Die erste wichtige Komponente, die wir modellierne müssen, sind die Datenstrukturen, die wir in der Datenbank speichern wollen.%
%
\item<2-> Es gibt dafür drei grundlegende Elemente für konzeptuelle Modelle: Entitäten, Attribute, und Entitätstypen.%
\end{itemize}%
%
\uncover<3->{%
\begin{definition*}[Entität]%
Eine \emph{Entität}~\inEN{entity} ist ein Objekt oder Ding das für sich selbst in der Welt existiert.\uncover<4->{ %
Sie kann von allen anderen Ojekten unterschieden werden.}%
\end{definition*}%
%
\uncover<5->{%
\begin{itemize}%
\item \DEZB\ ist der Student Herr Bibbo eine Entität\only<-5>{.}\uncover<6->{, das Module~\citetitle{programmingWithPython}\cite{programmingWithPython} ist eine Entität\only<-6>{.}\uncover<7->{, die Professorin Frau Bebba~教授 ist eine Entität\only<-7>{.}\uncover<8->{, und der Raum~\#36~202 ist auch eine Entität.}}}%
%
\item<9-> Entitäten können einfach in den Anforderungen oder Interviewnotizen entdeckt werden\only<-9>{.}\uncover<10->{ %
Sie entsprechen Eigennamen bzw.\ Hauptwörtern~\inEN{proper nouns}\cite{C1997ECAED}, also Wörtern, die eine bestimmte Sache identifizieren\cite{EOWM2025MWAMTD:CAPNWTDLWOGC}.}%
\end{itemize}}}%
\end{frame}%
%
\begin{frame}%
\frametitle{Attribute}%
\begin{definition*}[Attribut]%
\label{def:attribute}%
Ein \emph{Attribut}~\inEN{attribute}~$a$ ist charakterisiert by einem Namen und einer Domäne~\attrDomainb{a} mit Werten die es annehmen kann.%
\end{definition*}%
%
\uncover<2->{%
\begin{itemize}%
\item Attribute are sind Features und Characteristiken von Entitäten.%
%
\item<3-> Die Entitäten in unserem Modell werden durch die Werte ihrer Attribute definiert.%
%
\item<4-> Eine Studentin könnte durch ihren Namen, ihre ID, ihre Studenten-ID, ihre Mobiltelefonnummer, ihre Heimatadresse, und durch ihr \glslink{dateOfBirth}{Geburtsdatum} definiert werden.%
%
\item<5-> Ein Modul könnte durch seinen Titel, Syllabus, und Zusammenfassung definiert werden.%
%
\item<6-> Neben solchen Eigenschaften können \emph{Adjektive} in den Anforderungen oft als Attributwerte interpretiert werden\cite{C1997ECAED}, \DEzB\ rot, jung, erfolgreich, schwer, schnell\dots%
\end{itemize}%
}%
\end{frame}%
%
\begin{frame}[t]%
\frametitle{Domäne}%
%
\begin{definition*}[Domäne]%
Die \emph{Domäne}~\inEN{domain}~\attrDomainb{a} eines Attributs~$a$ ist die Menge der möglichen Werte, die es annehmen kann.%
\end{definition*}%
%
\uncover<2->{%
\begin{itemize}%
\item Wir können die Konzepte \emph{Domäne} und \emph{Datentype} im Kontext von konzeptuellen Modellen unterscheiden\cite{S2024D:LDMRMRA}.%
%
\item<3-> Ein Datentyp ist ein mathematisches Konzept wohingegen eine Domäne ein logisches Konzept ist.%
%
\only<-7>{%
\item<4-> \DEZB\ sind \sqlil{VARCHAR(100)}, \sqlilIdx{SMALLINT} und \sqlilIdx{REAL} alles Datentypen.%
}%
%
\only<-8>{%
\item<5-> Der Name eines Studenten ist hingegen ein logisches Konzept, auch wenn er durch einen Text-String dargestellt wird.%
}%
%
\item<6-> Die Mobiltelefonnummern oder ID-Nummern können auch als Strings dargestellt werden, sind aber auch bestimmte Zeichen und bestimmte Längen beschränkt.%
%
\item<7-> Ein \glslink{dateOfBirth}{Geburtsdatum} ist ein besonderes Datum, dessen Jahr aus einer bestimmten Menge möglicher Jahre kommt.%
%
\item<8-> Das Ergebnis eine Prüfung kann eine Ganzzahl zwischen 0 und 100 sein.%
%
\item<9-> Die Domäne eines Attributes ist also eine Kombination aus Datentyp und einer Semantik, die die gültigen Werte einschränkt.%
\end{itemize}}%
\end{frame}%
%
\begin{frame}%
\frametitle{Entitätstypen}%
%
\begin{definition*}[Entitätstyp]%
Die Menge aller Entitäten die die selben Attribute haben wird \emph{Entitätstype}~\inEN{entity type} genannt.%
\end{definition*}%
%
\uncover<2->{%
\begin{itemize}%
\item Während also die Entität Herr Bebbo ein einzelner Student ist, ist die Menge aller möglicher Studenten ein Entitätstyp.%
%
\item<3-> Während das Modul \citetitle{programmingWithPython}\cite{programmingWithPython} eine einzelne Entität ist, ist die Menge aller möglicher Module ein Entitätstyp.%
%
\item<4-> Die Professorin Frau Bebba~教授 ist eine einzelne Entität, aber die Menge aller möglichen Professoren repräsentiert einen Entitätstyp.%
%
\item<5-> Entitätstypen sind Appelative oder Gattungsnamen~\inEN{common nouns} die für Gruppen von Dingen stehen\cite{C1997ECAED,EOWM2025MWAMTD:CAPNWTDLWOGC}.%
\end{itemize}%
}%
\end{frame}%
%
\begin{frame}[t]%
\frametitle{Gute Praxis}%
%
\only<-4>{%
\begin{definition*}[Entitätstyp]%
Die Menge aller Entitäten die die selben Attribute haben wird \emph{Entitätstyp}~\inEN{entity type} genannt.%
\end{definition*}%
}%
%
\uncover<2->{%
\only<-6>{%
\begin{itemize}%
\item Beachten Sie den Plural \emph{Attribute} in der Definition von Entitätstyp!%
\end{itemize}%
}%
%
\uncover<3->{%
\only<-6>{%
\bestPractice{entityTypesHaveMultipleAttributes}{%
Nur Dinge mit mehreren Attributen sollten Entitätstypen sein.%
}}%
%
\uncover<4->{%
\begin{itemize}%
\item Es ergibt ja wenig Sinn, so ein Konzept wie \emph{Jahr} als Entitätstyp zu betrachten -- es hat ja nicht mehrere Attribute.
\end{itemize}%
%
\uncover<5->{%
\bestPractice{entityTypeSingleObject}{%
Jede Entität (und jeder Entitätstyp) soll genau ein Object (einen Objekttyp) aus der Realität modellieren (und nicht mehr als einen)\cite{SS2005EIDDDFDB:NQORD}.%
}%
%
\uncover<6->{%
\begin{itemize}%
\item Wenn wir \DEzB\ ein Autorennen modellieren, dann müssen Fahrer und Autos separate Entitätstypen sein, obwohl je ein Fahrer einem Auto zugeordnet ist.%
%
\item<7-> Wenn wir stattdessen den Namen des Fahrers und den Typ und die Eigenschaften der Autos in einen gemeinsamen Entitätstype modellieren, dann werden wir wahrscheinlich Redundanz erzeugen.%
%
\item<8-> Wenn mehrere Autos des selben Typs am Rennen teilnehmen, dann speichern wir deren Informationen mehrmals, jeweils mit einem anderen Fahrer.%
\end{itemize}%
}}}}}%
\end{frame}%
%
%
\begin{frame}%
\frametitle{Entitätenmengen}%
\begin{definition*}[Entitätenmenge]%
Eine \emph{Entitätenmenge}~\inEN{entity set} ist eine Untermenge eines Entitätstyps.\uncover<2->{ %
Sie ist eine Menge von einigen Entitäten eines Entitätstypes, die zu einer Zeit existieren.}%
\end{definition*}%
%
\uncover<3->{%
\begin{itemize}%
\item Herr Bebbo ist eine einzelne Studenten-Entität, die Menge aller möglichen Studenten bildet einen Entitätstyp, und die Studenten Herr Bebbo, Herr Bibbo und Herr Bobbo zusammen sind eine Entitätenmenge.%
%
\item<4-> Die Module \citetitle{programmingWithPython}\cite{programmingWithPython} und \citetitle{databases}\cite{databases} zusammen bilden eine Entitätenmenge, welche wiederum eine Untermenge des Entitätstyps fpr Module ist.%
%
\item<5-> Der mathematische Terminus \emph{Menge} ist hier wirklich korrekt\only<-5>{.}\uncover<6>{:}%
%
\item<6-> Alle Entitäten haben eine einzigartige Identität und können daher von allen anderen Objekten unterschieden werden.%
%
\item<7-> Es gibt keine zwei identischen Herr Bibbos.%
%
\item<8-> Daher können Studenten in eine Menge gruppiert werden und auch Entitätstypen sind Mengen.%
\end{itemize}}%
%
\end{frame}%
%
\begin{frame}%
\frametitle{Interessante Aspekte~1: Fenster zur Welt}%
\begin{itemize}%
\only<-10>{%
\item Wenn wir uns das konzeptuelle Design von Datenbanken auf diese Art anschauen, dann bemerken wir ein paar interessante Aspekte.%
}%
%
\only<-11>{%
\item<2-> Zum Beispiel modellieren wir immer nur ein kleines Fenster zur echten Welt.%
}%
%
\item<3-> Wenn wir über Studenten, Module, und Studiengänge als Entitätstypen sprechen, dann nur im Kontext unserer spezifischen Applikation.%
%
\item<4-> Natürlich gibt es in unserer realen, großen, weiten Welt Studenten auch an anderen Universitäten in anderen Ländern.%
%
\item<5-> Diese Studenten könnten andere Attribute haben als unsere.%
%
\item<6-> Sie haben wahrscheinlich keinen chinesische Ausweis~(中华人民共和国居民身份证).%
%
\item<7-> Sie spielen aber in unserem Modell von unserem kleinen Teil der Welt keine Rolle.%
%
\item<8-> Auch modellieren wir nur die Aspekte von Studenten, die für unsere Anwendung relevant sind.%
%
\item<9-> Wir modellieren weder deren Hobbies, Lieblingslied, Schuhgröße, Lieblingsessen noch deren Haarfarbe.%
%
\item<10-> Diese spielen nämlich auch keine Rolle in unserer Anwendung.%
%
\item<11-> Der Entitätentyp \inQuotes{Student} umfasst also nicht das komplette reale Konzept von Studenten.%
%
\item<12-> Er ist nur eine simplifizierte Abbildung dieses Konzepts in unsere Anwendung.%
\end{itemize}%
\end{frame}%
%
\begin{frame}%
\frametitle{Interessante Aspekte~2: OOP}%
%
\begin{itemize}%
\item Wenn Sie schon einen Kurs über Programmieren (wie \DEzB~\cite{programmingWithPython}) gehabt haben, dann wird sich diese Art zu modellieren etwas wie Objekt-orientierte Programmierung~(\glsShort{OOP}) anfühlen.%
%
\item<2-> Wir könnten uns Entitätstypen wir Klassen vorstellen und Entitäten wir deren Instanzen.%
%
\item<3-> Attribute wären dann, naja, Attribute.%
%
\item<4-> Ich denke, Datenbank-Theoretiker würden diese Art zu denken vielleicht nicht leiden, aber ich denke, dass sie nicht falsch ist.%
%
\item<5-> Das ist keine schlechte Analogie.%
%
\item<6-> Wenn wir dann später zu Beziehungen zwischen Entitäten kommen, passt sie vielleicht aber nicht mehr so gut.%
%
\end{itemize}%
\end{frame}%
%
\section{Entity-Relationship Diagramme~(ERDs)}%
%
\begin{frame}%
\frametitle{Spezifizieren}%
\begin{itemize}%
\item Wir wissen, dass Entitätstypen und ihre Attribute im Grunde Datenstrukturen beim Programmieren entsprechen.%
%
\item<2-> Sie sind eine der wichtigen Komponenten von konzeptuellen Modellen.%
%
\item<3-> Aber wie schreiben wir sie auf?%
%
\item<4-> Wie können wir sie spezifizieren?%
%
\end{itemize}%
\end{frame}%
%
\begin{frame}[t]%
\frametitle{Entity-Relationship Diagramme~(ERDs)}%
\begin{itemize}%
%
\item Es gibt eine grafische Methode zum Spezifizieren von Entitätstypen und ihren Beziehungen\only<-1>{.}\uncover<2->{: %
Entity-Relationship Diagramme~(\glsShort{ERD}s)\cite{C2002ERMHEFTALL,C1975TRMTAUVOD,C1976TERMTAUVOD,KW2012ASHOTEDAIM,WF1995DHQDM,B1990CMERMO}.}%
%
\item<3-> Es gibt eine Vielzahl von Notationen, die für \glsShort{ERD}s genutzt werden kann\only<-3,9->{.}\uncover<4->{:%
\begin{itemize}%
%
\item die ursprüngliche Noatationen von \citeauthor{B1969DSD}\cite{B1969DSD} und \citeauthor{C1975TRMTAUVOD}\cite{C1975TRMTAUVOD,C1976TERMTAUVOD} werden noch immer verwendet\only<-4>{.}\uncover<5->{,}%
%
\item<5-> die Krähenfußnotation~\inEN{Crow's Foot} wird sehr häufig eingesetzt\cite{E1976BDSMEWACE,CM2000MDMAUDA}\only<-5>{.}\uncover<6->{,}%
 %
\item<6-> es gibt auch die umfassende und standardisierte \glsShort{IDEF1X}~Syntax\cite{FIPSPUB184,ISOIECIEEE2012ITMLP2SASFII}\only<-6>{.}\uncover<7->{ und}%
%
\item<7-> auch \glsShort{UML}\cite{OMG2017OUMLOU,RMHOSMUUIIIIOPPTRS1997UNG,BRJ1999TUMLRM} wird verwendet.%
%
\end{itemize}%
}%
%
\item<8-> Es gibt viele Varianten von Diagrammen, die als \glsShort{ERD} verwendet werden können.%
%
\item<9-> \Citeauthor{S2024D:CDMERDE} listet neun verschiedene Varianten\cite{S2024D:CDMERDE}.%
%
\item<10-> \Citeauthor{SS2005EIDDDFDB:CDDRAAML} listet zwei grundlegende Methoden und gibt auch verschiedene Varianten von \glsShort{UML}-basierten Ansätzen an\cite{SS2005EIDDDFDB:CDDRAAML}.%
%
\item<11-> Die Notationen, die von \citeauthor{V1999C5DMS:CDUTERM} verwendet werden, sind wieder leicht anders\cite{V1999C5DMS:CDUTERM}.%
%
\item<12-> \glsShort{ERD}s werden also wahrscheinlich leicht anders aussehen, je nach dem, wer sie gemalt hat und welche Werkzeuge diese Person verwendet hat.%
%
\item<13-> Sie sind aber leicht zu verstehen, also sind diese Unterschiede nicht so wichtig.%
\end{itemize}%
\end{frame}%
%
%
\begin{frame}%
\frametitle{ERDs Zeichnen}%
\begin{itemize}%
%
\item Wir werden uns im Folgenden verschiedene \glsShort{ERD}s anschauen.%
%
\item<2-> Unser Ziel ist aber immer, \emph{anwendbares} Wissen zu lernen.%
%
\item<3-> Wir werden also nicht nur \glsShort{ERD}s anschauen, wir werden sie auch zeichnen.%
%
\item<4-> Das ist sehr einfach.%
%
\item<5-> Wenn Sie die grundlegende Syntax verstanden haben, dann können Sie diese Diagramme mit beliebiger Zeichensoftware malen\only<-5>{.}\uncover<6->{:%
%
\begin{itemize}%
\item Sie könnten Grafikprogramme wie \libreoffice\ \softwareStyle{Draw} oder \inkscape\cite{R2003DMEWI1APGTYJFBTPLVI,K2021TBOI} verwenden.\uncover<7->{ %
Dann müssen Sie aber alle Formen selber einzeln malen, was zu Inkonsistenzen führt und generell sehr aufwändig ist.}%
%
\item<8-> Dann gibt es Werkzeuge wie \pgmodeler\cite{AES2006PPDM} oder \mysqlWorkbench\cite{M2013MWDM}, die viel besser Unterstützung für \glsShort{ERD}s bieten.\uncover<9->{ %
Diese Werkzeuge sind aber an bestimmte Technologien gebunden~(hier: \postgresql\ und \mysql).\uncover<10->{ %
Sie sind daher für die Arbeit auf der konzeptuellen Ebene, die eben gerade technologieunabhängig seien soll, ungeeignet.}}%
%
\item<11-> Es gibt viele Werzeuge, die wir ausprobieren könnten {\dots} \citeauthor{B2025DS:7DMTC} listed über 70 in~\bracketCite{B2025DS:7DMTC}.%
%
\item<12-> Ich habe mich für \yEd\cite{SG2015MDAWY,Y2011YGEM} entschieden.\uncover<13->{ Das ist ein Graphen-Editor, der das Zeichnen und Arrangieren von \glsShort{ERD}s unterstützt und unabhängig von einem Datenmodell ist.}
\end{itemize}}%
\item<14-> Wir probieren diesen Editor an einem Beispiel aus, dann können Sie ihn selbst verwenden.%
\end{itemize}%
\end{frame}%
%
\section{ERDs mit yEd Zeichnen}%
%
\begin{frame}%
\frametitle{Vorbereitung}%
\begin{itemize}%
%
\item Kurze Notiz: Im Kontext von \glsShort{ERD}s werden Entitätstypen manchmal Entitäten genannt.%
%
\item<2-> Egal, lassen wir uns davon nicht verwirren.%%
%
\item<3-> Wir wollen nun den Entitätstype \emph{Student} modellieren.%
%
\item<4-> Von der Anforderungsanalyse wissen wir, dass jeder Student einen Namen hat, einen von der Regierung ausgestellten Ausweis~(ID, 中国公民身份号码\cite{GB116431999CIN}), und auch einen Studentenausweis, eine Mobiltelefonnummer, eine Adresse, und ein \glslink{dateOfBirth}{Geburtsdatum}~\inEN{date of birth, \glsShort{dateOfBirth}}.%%
%
\item<5-> Das wollen wir nun modellieren.%
\end{itemize}%
\end{frame}%
%
\begin{frame}[t]%
\frametitle{Student ERD mit yEd Zeichnen}%
\begin{itemize}%
%
\only<-2>{%
\item Wir öffnen \yEd.
}%
%
\only<-3>{%
\item<2-> Wir scrollen die \menu{Palette} rechts runter, bis wir \menu{Entity Relationship} finden.%
%
\item<3-> Wir klicken auf \menu{Entity Relationship}.
}%
%
\only<-4>{%
\item<4-> Wir klicken nun auf der \menu{Entity}-Symbol und ziehen es in das leere Dokument.%
}%
%
\only<-5>{%
\item<5-> Wir haben das \menu{Entity}-Symbol in das Dokument gezogen.%
}%
%
\only<-6>{%
\item<6-> Wir doppel-klicken in das neue Entity-Symbol im Dokument um seinen Namen zu bearbeiten.%
}%
%
\only<-7>{%
\item<7-> Wir ändern den Name des Entitätstyps in \inQuotes{Student} und drücken \keys{\return}.%
}%
%
\only<-9>{%
\item<8-> Der Name hat sich geändert.%
%
\item<9-> Wir klicken nun auf das \menu{Attribute}-Symbol in der Palette.%
}%
%
\only<-10>{%
\item<10-> Wir ziehen das Attributsymbol in unser Dokument.%
}%
%
\only<-11>{%
\item<11-> Wir doppel-klicken in es um seinen Namen zu ändern.%
}%
%
\only<-12>{%
\item<12-> Wir ändern den Name zu, nun ja, \inQuotes{Name}.%
}%
%
\only<-14>{%
\item<13-> Wir haben den Attributname geändert.%
%
\item<14-> Jetzt klicken wir auf das Verbindungs-Symbol in der Palette un ziehen es auf das Attribut.%
}%
%
\only<-15>{%
\item<15-> Wir lassen das Verbindungssymbol auf dem Attribut \emph{Name} fallen.%
}%
%
\only<-16>{%
\item<16-> Dann klicken wir auf den Entitätstyp und verbinden das Attribut \emph{Name} mit dem Entitätstyp \emph{Student}.%
}%
%
\only<-17>{%
\item<17-> Das Attribut \emph{Name} ist nun mit dem Entitätstyp \emph{Student} verbunden.%
}%
%
\only<-18>{%
\item<18-> Wir fügen das Attribut \emph{ID} auf genau die gleiche Art hinzu.%
}%
%
\only<-20>{%
\item<19-> Wir fügen noch mehrere weitere Attribute hinzu.%
%
\item<20-> Danach klicken wir auf das Menü \menu{\underline{F}ile}.%
}%
%
\only<-22>{%
\item<21-> Wir wollen das Dokument nun speichern.
%
\item<22-> Wir klicken auf \menu{S\underline{a}ve As}.%
}%
%
\only<-23>{%
\item<23-> Wir speichern das Diagramm unter dem Name \textil{erd_student_1} im \textil{graphml}-Format, in dem wir diesen Namen eingeben und dann auf \menu{Save} klicken.%
}%
%
\only<-24>{%
\item<24-> Wir können das Diagramm auch in ein Format exportieren, dass wir in anderen Dokumenten verwenden können.%
}%
%
\only<-25>{%
\item<25-> Dafür klicken wir auf \menu{\underline{E}xport} im Menü \menu{\underline{F}ile}.%
}%
%
\only<-26>{%
\item<26-> Wir wählen das \glsShort{SVG}-Format und klicken \menu{Save}.%
}%
%
\only<-27>{%
\item<27-> Wir lassen alle Einstellungen so wie sie sind und klicken \menu{OK}.%
}%
%
\item<28-> Das Ergebnis sieht genau so aus.%
\end{itemize}%
%
\locateGraphicTB{-3}{width=0.5\paperwidth}{graphics/yedErd/yedErdEntitiesA01scrollToErd}{0.25}{0.28}%
\locateGraphicTB{4}{width=0.5\paperwidth}{graphics/yedErd/yedErdEntitiesA02entity}{0.25}{0.28}%
\locateGraphicTB{5}{width=0.5\paperwidth}{graphics/yedErd/yedErdEntitiesA03dragEntity}{0.25}{0.28}%
\locateGraphicTB{6}{width=0.5\paperwidth}{graphics/yedErd/yedErdEntitiesA04doubleClickName}{0.25}{0.28}%
\locateGraphicTB{7}{width=0.5\paperwidth}{graphics/yedErd/yedErdEntitiesA05changeNameToStudent}{0.25}{0.28}%
\locateGraphicTB{8-9}{width=0.5\paperwidth}{graphics/yedErd/yedErdEntitiesA06nameIsStudentSelectAttribute}{0.25}{0.28}%
\locateGraphicTB{10}{width=0.5\paperwidth}{graphics/yedErd/yedErdEntitiesA07dragAttribute}{0.25}{0.28}%
\locateGraphicTB{11}{width=0.5\paperwidth}{graphics/yedErd/yedErdEntitiesA08doubleClickName}{0.25}{0.28}%
\locateGraphicTB{12}{width=0.5\paperwidth}{graphics/yedErd/yedErdEntitiesA09changeNameToName}{0.25}{0.28}%
\locateGraphicTB{13-14}{width=0.5\paperwidth}{graphics/yedErd/yedErdEntitiesA10changeNameChangedToNameSelectConnection}{0.25}{0.28}%
\locateGraphicTB{15}{width=0.5\paperwidth}{graphics/yedErd/yedErdEntitiesA11dragConnection}{0.25}{0.28}%
\locateGraphicTB{16}{width=0.5\paperwidth}{graphics/yedErd/yedErdEntitiesA12connectToStudent}{0.25}{0.28}%
\locateGraphicTB{17}{width=0.5\paperwidth}{graphics/yedErd/yedErdEntitiesA13connectedToStudent}{0.25}{0.28}%
\locateGraphicTB{18}{width=0.5\paperwidth}{graphics/yedErd/yedErdEntitiesA14addedID}{0.25}{0.28}%
\locateGraphicTB{19-20}{width=0.5\paperwidth}{graphics/yedErd/yedErdEntitiesA15addedAddrMobDoB}{0.25}{0.28}%
\locateGraphicTB{21-22}{width=0.5\paperwidth}{graphics/yedErd/yedErdEntitiesA16saveAs}{0.25}{0.28}%
\locateGraphicTB{23}{width=0.5\paperwidth}{graphics/yedErd/yedErdEntitiesA17saveAsErdStudent1}{0.25}{0.28}%
\locateGraphicTB{24-25}{width=0.5\paperwidth}{graphics/yedErd/yedErdEntitiesA18export}{0.25}{0.28}%
\locateGraphicTB{26}{width=0.5\paperwidth}{graphics/yedErd/yedErdEntitiesA19exportToSvg}{0.25}{0.28}%
\locateGraphicTB{27}{width=0.5\paperwidth}{graphics/yedErd/yedErdEntitiesA20exportToSvgDone}{0.25}{0.28}%
\locateGraphic{28}{width=0.5\paperwidth}{graphics/yedErd/yedErdEntitiesA21erd}{0.25}{0.33}%
\end{frame}%
%
\section{Mehr Attribute}
%
\begin{frame}%
\frametitle{Feedback}%
\begin{itemize}%
\item Wir nehmen also unser \glsShort{ERD} und diskutieren es mit unseren Partnern in der Universität.%
%
\item<2-> Wir wollen prüfen, ob das Model sinnvoll ist.%
%
\item<3-> Dabei entdecken wir ein paar Probleme\dots%
%
\end{itemize}%
\end{frame}%
%
\begin{frame}[t]%
\frametitle{Zusammengesetzte Attribute}%
%
\begin{itemize}%
\only<-7>{%
\item Namen sind komplizierter als wir dachten.%
%
\item<2-> Studenten haben einen Vorname, einen Familienname, und eine Anrede.%
%
\item<3-> Zum Beispiele ist der Name von \inQuotes{Herr Bebbo} eigentlich \inQuotes{Herr Fred Bebbo}\only<-3>{.}\uncover<4->{, wobei \inQuotes{Bebbo} der Familienname und \inQuotes{Fred} der Vorname sind.}%
}%
\item<4-> Um dem Rechnung zu tragen, sollte das \emph{Name}-Attribute kein einzelnes Attribut sein.%
%
\item<5-> Wir könnten es in ein \emph{zusammengesetztes Attribut} umbauen, dass aus einem Vorname~\inEN{given name}, Familienname~\inEN{family name} und einer Anrede~(engl.\ Salutation) besteht.%
%
\item<6-> Das System kann diese dann vernünftig zusammensetzen, wenn Studenten angesprochen oder Dokumente ausgestellt werden sollen.%
\end{itemize}%
%
\uncover<7->{%
\begin{definition*}[Zusammengesetztes Attribut]%
Ein \emph{zusammengesetztes Attribut}~\inEN{composite attribute} ist ein Attribut, das aus mehreren Teilen besteht, die eigene Namen und Domänen haben.
\end{definition*}%
\uncover<8->{%
\begin{definition*}[Einfaches Attribut]%
Ein \emph{einfaches Attribut}~\inEN{simple attribute} ist ein Attribut das nur einen einzelnen Namen und eine Domäne hat.\uncover<9->{ %
Einfache Attribute haben keine Komponenten.}
\end{definition*}%
}}
\end{frame}%
%
\begin{frame}[t]%
\frametitle{Probleme mit Namen}%
\begin{itemize}%
\only<-10>{%
\item Namen sind erstaunlich schwierig.%
}%
%
\only<-11>{%
\item<2-> In westlichen Kulturen wird der volle Name oft zusammengesetzt, in dem wir erst den Vorname und dann den Familiennamen schreiben.%
}%
%
\only<-12>{%
\item<3-> Deshalb ist der volle Name von Herrn Bebbo nämlich Fred Bebbo.%
}%
%
\only<-13>{%
\item<4-> In östlichen Kulturen ist es oft andersherum\only<-4>{.}\uncover<5->{:}%
}%
%
\only<-14>{%
\item<5-> Der Name des berühmten chinesischen Mathematikers 刘徽 wird als LIU Hui im lateinischen Alphabet transkribiert\cite{OR2003LH,S1998LHATFGAOCM,Y2024COACMMLHFHTIOMACE}.%
}%
%
\only<-15>{%
\item<6-> Dabei ist LIU~(刘) der Familienname und kommt zuerst und der Vorname Hui~(徽) kommt danach.%
}%
%
\only<-16>{%
\item<7-> Abhängig vom kulturellen Hintergrund könnte das System also den Namen korrekt zusammensetzen.%
}%
%
\only<-17>{%
\item<8-> Das klingt aber nach einem grässlichen Problem und nach extrem fehleranfälligem Kode.%
}%
%
\only<-18>{%
\item<9-> Wir entscheiden uns also dagegen, den Namen aus Vor- und Nachname zusammenzusetzen.%
}%
%
\only<-19>{%
\item<10-> Stattdessen benutzen wir zwei andere Komponenten des Namensattributes\only<-10>{.}\uncover<11->{:%
\begin{enumerate}%
\item Der volle Name~\inEN{Full Name} ist der komplette offizielle Name der Person, genauso geschrieben wie im Ausweis.\uncover<12->{ %
Das wäre so etwas wie 刘徽 für eine chinesische Person und Fred Bebbo für den Westerler Herr Bebbo.}%
\item<13-> Als zweite Komponente speichern wir die volle Anrede~\inEN{Salutation}.\uncover<14->{ %
Für 刘徽 könnte das 刘教授 sein und für Fred Bebbo vielleicht Herr Bebbo.}%
\end{enumerate}%
}%
%
\item<15-> Vielleicht erlauben wir den Studenten später, die \emph{Salutation} selbst auszuwählen.%
%
\item<16-> Der \emph{Full Name} wird vom Administrator genauso eingegeben, wie er auf dem Ausweis steht.%
}%
%
\item<17-> Unser System würde die \emph{Salutation} verwenden, wenn es Nachrichten an Herrn Bebbo schickt.%
%
\item<18-> Auf Zertifikaten wird stattdessen der \emph{Full Name} verwendet.%
%
\item<19-> Ein Zusammengesetztes Attribut wird in \glsShort{ERD}s als ein Attribut gemalt, dass mit anderen Attributen verbunden ist.%
\end{itemize}%
\locateGraphic{20}{width=0.7\paperwidth}{graphics/erds/erdStudent2a}{0.15}{0.46}%
\end{frame}%
%
\begin{frame}[t]%
\frametitle{Mehrwertige und Optionale Attribute}%
%
\only<-5>{%
\begin{definition*}[Mehrwertiges Attribut]%
Eine Entität kann mehrere Werte für eine mehrwertiges~\inEN{multivalued} Attribut haben.%
\end{definition*}%
}%
%
\uncover<2->{%
\only<-9>{%
\begin{definition*}[Einwertiges Attribut]%
Eine Entität hat (höchstens) einen Wert für ein einwertiges~(\inEN{single-valued}) Attribut.%
\end{definition*}%
}%
\uncover<3->{%
\begin{definition*}[Optionales Attribut]%
Eine Entität kann einen Wert für ein optionales~\inEN{optional} Attribut haben, oder auch nicht.\uncover<4->{ %
Wenn ein Attribut keinen Wert hat, dann wird das oft durch \sqlilIdx{NULL} dargestellt.}%
\end{definition*}%
%
\uncover<5->{%
\begin{itemize}%
\item Ein optionales einwertiges Attribut hat entweder einen Wert oder keinen Wert.%
\item<6-> Ein nicht-optionales einwertiges Attribut hat immer genau einen Wert.%
\item<7-> Ein optionales mehrwertiges Attribut hat entweder keinen Wert, einen Wert, oder mehrere Werte.%
\item<8-> Ein nicht-optionales mehrwertiges Attribut hat entweder einen Wert oder mehrere Werte.%
\item<9-> Ein mehrwertiges Attribut wird in \glsShort{ERD} durch zwei ineinanderliegende Ellipsen dargstellt.%
\item<10-> Optionale Attribute werden durch kleine kreise am Ende der Linie die sie mit den Entitätstypen verbinden dargestellt.%
\end{itemize}%
}}}%
\end{frame}%
%
\begin{frame}%
\frametitle{Mobiltelefonnummern}%
\begin{itemize}%
\item Wir stellen fest, dass ein Student mehrere Mobiltelefonnummern haben kann.%
%
\item<2-> Das kann mit einem mehrwertigen Attribut modelliert werden.%
%
\item<3-> Bei genauerer Betrachtung stellen wir fest, dass das Mobiltelefonnummer-Attribut am besten als optionales mehrwertiges Attribut modelliert werden sollten.%
%
\item<4-> Es ist heutzutage zwar eher unwahrscheinlich, dass ein Student keine Mobiltelefonnummer hat, aber es könnte sein.%
%
\item<5-> Vielleicht trifft das ja auf ausländische Austauschstudenten (留学生) zu, die gerade erst nach China gekommen sind.%
%
\end{itemize}%
\end{frame}%
%
\begin{frame}%
\frametitle{Ausweisnummern}%
\begin{itemize}%
\item Und wo wir gerade bei dem Thema Austauschstudenten sind\dots%
%
\item<2-> They probably do not have a Chinese~ID~(中国公民身份号码\cite{GB116431999CIN}).
%
\item<3-> Deshalb sollte auch ID ein optionales Attribut sein.
%
\end{itemize}%
\end{frame}%
%
\begin{frame}[t]%
\frametitle{Neues ERD}
Wir bekommen das folgende neue \glsShort{ERD}.%
%
\locateGraphic{}{width=0.6\paperwidth}{graphics/erds/erdStudent2}{0.2}{0.18}%
\end{frame}%
%
\begin{frame}[t]%
\frametitle{Adressen}%
\begin{itemize}%
\only<-7>{%
\item Wir benerken auch, dass \emph{Address} wahrscheinlich auch ein zusammengesetztes Attribut seien sollte.%
\item<2-> Eine Adresse ist ja nicht nur eine Zeile Text.%
\item<3-> Sie besteht aus Land, Provinz, Stadt, Distrikt, Postleitzahl, Straße und Hausnummer.%
\item<4-> Das zu modellieren wäre umständlich.%
\item<5-> Wir entscheiden uns dafür, dass alles nach der Postleitzahl~\inEN{postal code} einfach eine Zeichenkette mit Namen \emph{street address} werden soll.%
\item<6-> Ein automatisches kann Informationen die feiner auflöst als eine Postleitzahl wahrscheinlich sowieso nicht sinnvoll verarbeiten.%
%
\item<7-> Während wir diese Adressverarbeitung mit unseren Partnern diskutieren, stellen wir fest, dass Studenten auch mehrere Adressen haben können.%
}%
\item<8-> Wir bekommen also ein neues \glsShort{ERD}.%
\end{itemize}%
\locateGraphic{8}{width=0.8\paperwidth}{graphics/erds/erdStudent3}{0.1}{0.18}%
\end{frame}%
%
%
\begin{frame}[t]%
\frametitle{Abgeleitete Attribute}%
%
\only<-7>{%
\begin{definition*}[Abgeleitetes Attribut]%
Ein abgeleitetes Attribut~\inEN{derived attribute} wird aus den Werten anderer Attribute berechnet.%
\end{definition*}%
}%
%
\uncover<2->{%
\begin{itemize}%
\only<-7>{%
\item Abgeleitete Attribute müssen nicht in der Datenbank gespeichert werden.%
%
\item<3-> Ein typisches Beispiel ist das Alter einer Person.%
%
\item<4-> Wenn wir den Geburtstag~(\glsShort{dateOfBirth}) und das aktuelle Datum kennen, dann kann das Alter berechnet werden.%
%
\item<5-> Das geht sehr schnell.%
%
\item<6-> Es gibt keinen Grund, das Alter~\inEN{age} einer Person in der Datenbank zu speichern.%
}%
%
\item<7-> Abgeleitete Attribute werden in \glslink{ERD}{ERDs} mit Ellipsen mit gestrichelten Linien dargestellt.%
%
\end{itemize}}%
\locateGraphic{8}{width=0.8\paperwidth}{graphics/erds/erdStudent4}{0.1}{0.18}%
\end{frame}%
%
\section{Zusammenfassung}%
%
\begin{frame}%
\frametitle{Zusammenfassung: Entitäten}%
\begin{itemize}%
\item Wir haben nun sehr viel über Entitäten und Attribute gelernt.%
%
\item<2-> Eine Entität modelliert ein Objekt aus der realen Wert, wie \DEzB\ eine Person, ein Ding, einen Ort, ein Ereignis, oder ein Konzept.%
%
\item<3-> Eine Entität existiert nicht einfach so.%
%
\item<4-> Sie ist ein Tupel aus ihren Attributwerden.%
%
\item<5-> Während Entitäten einzelne Objekte beschreiben, beschreiben Entitätstypen Gruppen von Objekten, wie Personen, Dinge, Orte, eine Klasse von Ereignissen, usw.%
\end{itemize}%
\end{frame}%
%
\begin{frame}%
\frametitle{Zusammenfassung: Attribute}%
\begin{itemize}%
\only<-11>{%
\item Attribute können ebenfalls verschiedene Eigenschaften haben.%
}%
%
\only<-12>{%
\item<2-> Sie können einfach sein, also einen Wert haben, der nicht weiter unterteilt werden kann\only<-2>{.}\uncover<3->{, wie eine Telefonnummer.}%
}%
%
\only<-13>{%
\item<4-> Sie können zusammengesetzt sein, also Teilen bestehen\only<-4>{.}\uncover<5->{, wie Adressen, die aus Land, Provinz, Stadt, usw.\ bestehen.}%
}%
%
\only<-14>{%
\item<5-> Es ist nicht immer klar, wie ein Attribut modelliert werden soll.%
}%
%
\item<6-> Zum Beispiel haben unsere Studenten ein Attribut \glsFull{dateOfBirth}.%
%%
\item<7-> Technisch haben wir bereits gelernt, dass ein Datum ein atomarer Datentyp in \glsShort{SQL} ist.%
%
\item<8-> Wir können \glsShort{dateOfBirth} also durchaus als einfaches Attribut modellieren.%
%
\item<9-> Wir könnten es aber auch als zusammengesetztes Attribut aus Jahr, Monat, und Tag modellieren.%
%
\item<10-> Das würde allerdings kein vernünftig denkender Mensch machen, denn dann muss man ja auch Dinge wie Zeitzonen und verschiedene Kalender beachten\dots%
%
\item<11-> Eine Adresse kann ein einfach Textstring oder ein zusammengesetztes Attribut sein.%
%
\item<12-> Die Entscheidung, wie wir Attribute modellieren, hängt oft von den Daten ab, die wir brauchen.%
%
\item<13-> Wenn wir das Geburtsdatum als einfaches Attribut vom Datentyp \inQuotes{Datum} speichern, dann kann wir sehr einfach das Geburtsjahr oder das Alter einer Person berechnen.%
%
\item<14-> Speichern wir die Adresse jedoch als einfaches Attribut, dann wird es eben sehr schwer, das Land, die Provinz, oder die Stadt automatisch zu extrahieren.%
%
\item<15-> Attribute können zusätzlich noch einwertig oder mehrwertig sein, optional oder nicht-optional.%
\end{itemize}%
\end{frame}%
%
\begin{frame}%
\frametitle{Zusammenfassung}%
\begin{itemize}%
\item Wir können nun die Daten, die wir in unsere Datenbank speichern, schon ziemlich gut modellieren.%
%
\item<2-> In einer Datenbank geht es aber nicht nur um Daten.%
%
\item<3-> Es geht auch um Einschränken der Daten und die Beziehungen zwischen Daten.%
%
\item<4-> Und als nächstest lernen wir eine bestimmte Einschränkung kennen\dots%
\end{itemize}%
\end{frame}%
%
\endPresentation%
\end{document}%%
\endinput%
%
